% !Mode:: "TeX:UTF-8"
\chapter{实验结果与分析}
\section{实验环境及配置}
% 硬件平台（GPU、机械臂型号、手爪参数）。
% 软件环境与评价指标（Precision, Recall, F1, Success Rate, 抓取分布统计指标）。

\section{点云特征检测实验分析}
\subsection{对比实验}
% 与 PIE-Net, PCEDNet, BoundED 等 SOTA 方法在 ABC-Corner 数据集上的对比（重点展示角点检测的优越性）。
\subsection{消融实验}
% 特征维度选择的影响、邻域尺度数量的影响、过采样策略的有效性。
\subsection{结果可视化}
% 展示不同物体点云的边缘与角点分割效果。

\section{抓取规划仿真实验分析}
\subsection{抓取分布多样性评估}
% 引入标准差(STD)、极差(R)等指标，对比本文方法与 GraspIt!, GPD, WRS 的抓取空间覆盖能力。
\subsection{典型物体抓取规划结果}
% 针对四面体、不规则棱柱、日用品等物体的规划结果展示（突出六种接触组合）。
\subsection{参数敏感性分析}
% 采样体素大小、检测器角度阈值等参数对规划结果的影响。

\section{真实机器人抓取实验}
% 实验设置与流程（物体放置、运动规划、抓取执行）。
% 多类别物体（工业件、日用品、异形件）的抓取成功率统计。
% 对边缘和角点接触抓取的实际效果验证。
\section{本章小结}
